% =====================================================================
%  Chapter 7 — Ultralight Bosons
% =====================================================================

\chapter{Ultralight Bosons and Wave Dark Matter}
\label{ch:ultralight}

One candidate class is not a particle in the usual sense. A field of very
small mass behaves as a coherent quantum wave over galactic distances, and
the dark matter in a galaxy is then a single quantum state rather than a
gas of particles. The phenomenology that follows is distinctive, and the
experimental consequences are different in kind from everything else in this
book.

\section{The mechanism and the mass scale}A field of mass $m$ has a de Broglie wavelength
$\lambda_{\rm dB} = h/(mv)$, where $v$ is the virial velocity in a halo. If
this wavelength is comparable to or larger than the halo's core, the field
cannot be described as a collection of particles at all: interference
effects suppress the formation of small structure.

The mass scale where this happens in a Milky Way halo is
$m \sim 10^{-22}\ \eV$, and it is a coincidence of no deep significance that
this is close to the mass scale on which the de Broglie wavelength equals
the size of a dwarf galaxy. The same scale also produces solitons in halo
centres, as Chapter~\ref{ch:smallscale} described, and suppresses the
satellite population.

The three effects are not independent. A field light enough to make large
solitons also erases the smallest structure, and the CMB lensing constraint
on small-scale power limits how light the field can be.

\section{What the field predicts and what it costs}The prediction set is specific and testable, which is
the strongest thing that can be said for the class.

The \emph{core structure} prediction is the soliton of Chapter~\ref{ch:smallscale}.
The \emph{abundance prediction} is a suppression of the satellite population
below a characteristic mass, and a modification of the low-mass end of the
halo mass function. The \emph{lensing prediction} is a characteristic
scale-dependent phase shift in the gravitational potential, which changes the
relation between the total mass and the shear in a wave-agnostic way: a
fixed amount of mass produces less shear in a wave-like field than in a
particle-like one, at the same scale. This effect has been measured and
constrained for wavelike dark matter in galaxy clusters and in the strong
lensing of individual systems.

The costs are the standard ones: a boson has not been observed, the
production mechanism must be identified, and --- decisively --- the same
wavelength that produces the phenomenology constrains it cosmologically.
Chapter~\ref{ch:smallscale} covered the soliton-host tension and the 2026
CMB lensing constraint. Both are currently unfavourable to the lightest
masses.

\section{Decoupling the class into three}The phrase ``ultralight dark matter'' covers three
distinct objects, and the field's habit of treating them as one has cost it
clarity.

The \emph{axion-like particle} is the pseudo-Nambu--Goldstone boson of a
broken global symmetry, with a coupling to electromagnetism that is
suppressed by the symmetry-breaking scale. The QCD axion is the special
case in which the symmetry is the Peccei--Quinn one, and it is the only
member of the class with a completely specified production mechanism. Other
ALPs exist and are better motivated as dark matter in some scenarios, because
their mass--coupling relation is not fixed by the strong CP problem.

The \emph{fuzzy dark matter} term was coined for the $\sim 10^{-22}$~eV
case specifically, where the soliton and the satellite suppression are the
relevant phenomenology.

The \emph{axion-like particle} mass range of interest for the wave
phenomenology extends from about $10^{-24}$ to $10^{-19}$~eV, and the
overlap with the QCD axion band is empty: the QCD axion, in the
post-inflationary scenario, is far heavier than the wave-phenomenology
range.

\section{The production problem}For a field with a mass in the wave-phenomenology range, the
misalignment mechanism that produces the QCD axion does not work. The
condition for misalignment to be the dominant production route is that the
initial misalignment angle be of order one, which requires that the field's
energy density at the beginning of the field-dominated era be no larger than
the total energy density. For a very light field, the initial condition
problem is severe, and the misalignment angle must be extremely finely
tuned --- or the field must be produced by a different mechanism entirely.

This is the point at which the ultralight class runs into a problem that the
WIMP and the axion do not have, and it is a problem about origins, not
about phenomenology. It is the subject of Chapter~\ref{ch:production}, and it
is one of the strongest arguments for the view that the origin of the relic
density is downstream of discovery rather than upstream of it.

\section{Decision criteria}An ultralight boson would be confirmed by a
measurement of the phase shift in gravitational lensing at the predicted
scale dependence, or by the observation of a soliton profile in a halo where
the alternatives are excluded, or --- the decisive case --- by the
laboratory detection of the field itself through its electromagnetic
coupling. It would be excluded, in the lightest-mass region, by continued
CMB lensing and large-scale-structure observations that show small-scale
power where the model requires its absence.

\begin{keytakeaways}
\item Wave-like dark matter predicts a coherent set of effects --- solitons,
satellite suppression, scale-dependent lensing phase --- that are testable
and largely independent of each other.
\item The mass range is split into three distinct objects with different
production mechanisms, and the literature's habit of calling all of them
ultralight dark matter obscures the actual constraints.
\item The production of a very light field from plausible initial conditions
is a genuine unsolved problem, and it is an argument for discovery-first
strategy.
\item The lightest masses are currently under pressure from CMB lensing.
\end{keytakeaways}
